% ======================================================================
% Ontologie der Kontinua -- Core 1.3.3
% Cross-K-Modul: crossk_k10_k11.tex
% Cross-Level-Relationen zwischen K10 und K11
% ======================================================================

\subsubsection{K10–K11-Querverbindung}
\label{sec:crossk-k10-k11}

Der Übergang $K_{10} \rightarrow K_{11}$ markiert die Entstehung des
\textbf{trans-metatheoretischen Kontinuums} aus dem metatheoretischen
Kontinuum. Während $K_{10}$ funktoriale Abbildungen, Meta-Logik und
kategorielle Strukturen bereitstellt, die theoretische Systeme miteinander verbinden,
führt $K_{11}$ ein:
\begin{itemize}
  \item universelle Operatorklassen, die über mehrere Meta-Rahmenwerke wirken,
  \item höherordentliche Verträglichkeitsbedingungen für Meta-Modelle,
  \item globale Schwellen für trans-meta Kohärenz,
  \item Flüsse, welche ganze Familien von Meta-Modellen restrukturieren,
  \item Zyklen, die Meta-Ebene selbst regulieren,
  \item eine Einschränkungsschicht, die die Modellentwicklung über alle Stufen
        $K_0\dots K_{10}$ konsistent innerhalb eines höheren Umschlags hält.
\end{itemize}

In Core~1.3.3 wird $K_{11}$ als aktive obere Doktrin behandelt, nicht als
unaufgelöster Platzhalter. Dieses Kapitel fasst die strukturelle Rolle innerhalb
der geförderten Cross-Level-Architektur zusammen.

% ----------------------------------------------------------------------
\subsubsection{1. Stufen und ihre Rollen}

\paragraph{K10 (metatheoretisches Kontinuum).}
$K_{10}$ liefert:
\begin{itemize}
  \item Achsen $A^{10}$ (Funktoren, Meta-Logik, Kompatibilitätskategorien),
  \item Potentiale $P^{10}$ (Meta-Kohärenz, Ausdrucksstärke, Universalität),
  \item Schwellen $\Theta^{10}$ (Diagrammkonsistenz,
        formale Meta-Logik-Korrektheit),
  \item Flüsse $J^{10}$ (Interpretation, Übersetzung, Rekonstruktion),
  \item Zyklen $C^{10}$ (Meta-Stabilität, Verfeinerungsschleifen).
\end{itemize}

\paragraph{K11 (trans-metatheoretisches Kontinuum).}
$K_{11}$, in der Core~1.3.3-Formulierung, führt ein:
\begin{itemize}
  \item Achsen $A^{11}$ zur Beschreibung von Relationen zwischen vollständigen
        Klassen von Meta-Rahmenwerken,
  \item Potentiale $P^{11}$ zur Quantifizierung globaler Kohärenz über Meta-Ebenen,
  \item Schwellen $\Theta^{11}$ für trans-meta Konsistenz und Kompatibilität,
  \item Flüsse $J^{11}$, die Familien von Metamodellen restrukturieren,
  \item Zyklen $C^{11}$ zur Regulierung der Evolution von $K_{10}$-Strukturen,
  \item Maß $k(K_{11})$ für universelle strukturelle Kohärenz.
\end{itemize}

$K_{11}$ dient als Einschränkungsschicht für alle unteren Modellierungsstufen.

% ----------------------------------------------------------------------
\subsubsection{2. Vererbte und neue Achsen sowie Schwellen}

\paragraph{Vererbte Struktur.}
Metamodelle in $K_{10}$ werden zu Objekten; $K_{11}$ bringt Morphismen zwischen
Metamorphismen:
\[
\mathsf{MetaFunctors}(K_{10}) \rightarrow \mathsf{TransFunctors}(K_{11}).
\]

\paragraph{Neue Achsen.}
Achsen in $A^{11}$ umfassen:
\begin{itemize}
  \item $A_{\mathrm{ultra}}$ — Operatorklassen auf Kategorien von Meta-Modellen,
  \item $A_{\mathrm{compat}}^{(11)}$ — trans-meta Kompatibilitätsrelationen,
  \item $A_{\mathrm{reflect}}^{(11)}$ — Achsen höherer Selbstreflexion,
  \item $A_{\mathrm{limit}}$ — Einschränkungen aus universeller Konsistenz.
\end{itemize}

\paragraph{Neue Schwellen.}
$\Theta^{11}$ umfasst:
\begin{itemize}
  \item \textbf{trans-meta Kohärenzschwelle} — Versagen führt zum Kollaps von $K_{11}$,
  \item \textbf{Kompatibilitätsschwelle} — inkompatible Meta-Rahmenwerke können
        nicht koexistieren,
  \item \textbf{universelle Konsistenzschwelle} — vermeidet Paradoxien auf
        trans-meta Ebene,
  \item \textbf{Ausdrucksstärkeschwelle} — minimale Repräsentationssuffizienz über
        $K_0\dots K_{10}$.
\end{itemize}

Verletzung von $\Theta^{10}$ verhindert das Auftreten von $K_{11}$.

% ----------------------------------------------------------------------
\subsubsection{3. Cross-Level-Flüsse, Zyklen und Spannungen}

\paragraph{Flüsse.}
Flüsse in $K_{11}$ erweitern Meta-Flüsse:
\[
J^{11} = (J_{\mathrm{ultra}},\ J_{\mathrm{crossmeta}},\
         J_{\mathrm{lift2}},\ J_{\mathrm{constraint}}),
\]
wobei:
\begin{itemize}
  \item $J_{\mathrm{ultra}}$ — Restrukturierung von Kategorien von Metamodellen,
  \item $J_{\mathrm{crossmeta}}$ — Übersetzung zwischen inkompatiblen Meta-Ebenen,
  \item $J_{\mathrm{lift2}}$ — Heben von Meta-Frameworks in $K_{11}$,
  \item $J_{\mathrm{constraint}}$ — Durchsetzung globaler struktureller Grenzen.
\end{itemize}

\paragraph{Zyklen.}
$C^{11}$ umfasst höhergradige Zyklen wie:
\[
C_{\mathrm{ultra}}:\
\text{Metamodell} \rightarrow \text{Lift} \rightarrow \text{Constraint}
\rightarrow \text{Verfeinerung} \rightarrow \text{Metamodell}.
\]

Diese Zyklen regulieren die Evolution des Meta-Kontinuums.

\paragraph{Spannung.}
Cross-Level-Spannung:
\[
T_{10,11} =
f(\Theta^{11} - P^{11},\ A^{11} - A^{10},\
  J^{11} - J^{10},\ \Theta^{10} - \Theta^{11}).
\]

Übermäßige Spannung kollabiert $K_{11}$ zu einem einfacheren metatheoretischen
Regime.

% ----------------------------------------------------------------------
\subsubsection{4. Geburts- und Todesbedingungen über die Stufen}

\paragraph{Geburt von $K_{11}$.}
$K_{11}$ entsteht, wenn:
\[
T_{10} > \Theta_{10,\mathrm{dim}}
\quad\text{und}\quad
A^{11} \in M_{11} \setminus A(K_{10}),
\]
d.\,h. wenn Unterscheidungen zwischen vollständigen Meta-Framework-Klassen
erforderlich werden.

Erweiterung:
\[
\Omega(K_{11}) = \Omega(K_{10}) \cup \Delta\Omega_{\mathrm{ultra}}.
\]

\paragraph{Todesausbreitung.}
Wenn $\Theta^{10}$ versagt, kann die trans-meta Ebene nicht bestehen.

Wenn $\Theta^{11}$ versagt (Inkompatibilität, meta-paradoxer Konflikt), kollabiert
$K_{11}$, während $K_{10}$ bestehen bleibt:
\[
\Omega(K_{11}) \rightarrow \varnothing.
\]

% ----------------------------------------------------------------------
\subsubsection{5. Konzeptionelle Beispiele}

\paragraph{Inkompatible Meta-Rahmenwerke vermitteln.}
$K_{11}$ kann einen höherstufigen Kompatibilitätsoperator liefern, der
Meta-Frameworks verknüpft, die sich in $K_{10}$ nicht direkt abbilden.

\paragraph{Universelle Einschränkungen.}
$K_{11}$ erzwingt globale Konsistenz über alle Modellierungsstufen:
\[
K_{11} : \{K_0,\dots,K_{10}\} \rightarrow \text{kohärente Hülle}.
\]

\paragraph{Kollapsfall.}
Entsteht ein trans-meta Inkonsistenz:
\[
T_{10,11} > \Theta^{11}_{\mathrm{coh}},
\]
fällt $K_{11}$ zurück in $K_{10}$.

So lässt sich der K10→K11-Übergang in begrenzter Form als Entstehen
universaler Operatoren und Nebenbedingungen zusammenfassen, die der Beziehung und
Regulierung ganzer Klassen metatheoretischer Rahmenwerke dienen.
